\section{Part VI: The Register — Retracted, Corrected, and Audited}

\subsection{Color}

A degenerate block in a scale pattern — a set of directions carrying the same local unit — has no preferred labelling within it.  When a defect sits in such a block and a circuit is made around it, the scale winds and the block labels may be permuted.  Asking when such a configuration is globally consistent (when the labelling closes up after the loop) yields three striking results: defects with genuine permutations cannot exist alone (confinement), bound states must contain exactly as many constituents as the block size, and charge comes in units of $1/n$ for a block of size $n$.  These theorems are mathematically correct but the identification with colour has been withdrawn — the geometry of the scale gives a block of size two, not three, as registered in \texttt{ColorAudit.lean}.

\begin{definition}
The \textbf{monodromy} of a defect in a degenerate block is a pair: a permutation of the block labels and an integer winding per direction.  Going once around a defect, the scale winds (forced by discrete scale invariance) and the block labels may be permuted (forced by the absence of preferred labelling).
\begin{align*}
\text{BlockMonodromy}(n) = \{\text{perm} : \text{Equiv.Perm}(\text{Fin } n), \text{ wind} : \text{Fin } n \to \mathbb{Z}\}
\end{align*}
\lean{Color.BlockMonodromy}
\end{definition}

\begin{definition}
The composition of two monodromies combines their permutations multiplicatively and their windings additively:
\begin{align*}
M \circ N &: \text{perm} = M.\text{perm} \cdot N.\text{perm}, \quad \text{wind} = \lambda a. M.\text{wind}(a) + N.\text{wind}(a)
\end{align*}
\lean{Color.BlockMonodromy.comp}
\end{definition}

\begin{definition}
The opposite defect, obtained by reversing the circuit direction, has inverted permutation and negated winding:
\begin{align*}
M^{-1} &: \text{perm} = M.\text{perm}^{-1}, \quad \text{wind} = -M.\text{wind}
\end{align*}
\lean{Color.BlockMonodromy.anti}
\end{definition}

\begin{definition}
The vacuum state is the trivial monodromy — identity permutation and zero winding everywhere.
\begin{align*}
\text{triv}(n) &: \text{perm} = 1, \quad \text{wind}(a) = 0
\end{align*}
\lean{Color.BlockMonodromy.triv}
\end{definition}

\begin{definition}
The total scale winding — the charge — is the sum of windings over all directions:
\begin{align*}
\text{charge}(M) = \sum_{a=0}^{n-1} M.\text{wind}(a)
\end{align*}
\lean{Color.BlockMonodromy.charge}
\end{definition}

\begin{proposition}
Composition acts on the permutation part by multiplication:
\begin{align*}
(M \circ N).\text{perm} = M.\text{perm} \cdot N.\text{perm}
\end{align*}
\end{proposition}
\begin{proof}
Immediate from the definition of $\circ$.
\end{proof}
\lean{Color.BlockMonodromy.comp\_perm}

\begin{proposition}
The opposite monodromy inverts the permutation part:
\begin{align*}
M^{-1}.\text{perm} = M.\text{perm}^{-1}
\end{align*}
\end{proposition}
\begin{proof}
Immediate from the definition of the opposite monodromy.
\end{proof}
\lean{Color.BlockMonodromy.anti\_perm}

\begin{proposition}
The vacuum monodromy has trivial permutation: $\text{triv}(n).\text{perm} = 1$.
\end{proposition}
\begin{proof}
Immediate from the definition of $\text{triv}$.
\end{proof}
\lean{Color.BlockMonodromy.triv\_perm}

\begin{proposition}
Composing two monodromies gives the sum of charges.
\begin{align*}
\text{charge}(M \circ N) = \text{charge}(M) + \text{charge}(N)
\end{align*}
\begin{proof}
Both sides are sums over the finite index set of directions: $\text{charge}(M\circ N)=\sum_a (M\circ N).\text{wind}(a) = \sum_a (M.\text{wind}(a)+N.\text{wind}(a))$ by definition of composition, and this splits by additivity of the finite sum into $\sum_a M.\text{wind}(a) + \sum_a N.\text{wind}(a) = \text{charge}(M)+\text{charge}(N)$.
\end{proof}
\lean{Color.BlockMonodromy.charge\_comp}
\end{proposition}

\begin{proposition}
The opposite defect has opposite charge.
\begin{align*}
\text{charge}(M^{-1}) = -\text{charge}(M)
\end{align*}
\begin{proof}
$\text{charge}(M^{-1}) = \sum_a M^{-1}.\text{wind}(a) = \sum_a (-M.\text{wind}(a))$ by definition of the opposite monodromy, and a finite sum of negatives is the negative of the sum, giving $-\sum_a M.\text{wind}(a) = -\text{charge}(M)$.
\end{proof}
\lean{Color.BlockMonodromy.charge\_anti}
\end{proposition}

\begin{definition}
A configuration is \textbf{observable} when its block labelling closes up after the circuit — the permutation is the identity.
\begin{align*}
\text{Observable}(M) :\Leftrightarrow M.\text{perm} = 1
\end{align*}
\lean{Color.BlockMonodromy.Observable}
\end{definition}

\begin{proposition}
The vacuum monodromy is observable: $\text{Observable}(\text{triv}(n))$.
\end{proposition}
\begin{proof}
The vacuum has trivial permutation, which is the identity by definition of
observability.
\end{proof}
\lean{Color.BlockMonodromy.triv\_observable}

\begin{theorem}[Confinement]
A defect whose monodromy permutes the block genuinely cannot exist on its own.  This is not a force growing with separation; the isolated configuration is simply not single-valued.
\begin{align*}
\text{If } M.\text{perm} \neq 1, \text{ then } \neg \text{Observable}(M)
\end{align*}
\begin{proof}
Immediate from the definitions: $\text{Observable}(M)$ unfolds to $M.\text{perm}=1$, so the hypothesis $M.\text{perm}\neq1$ is literally the negation of the conclusion.
\end{proof}
\lean{Color.BlockMonodromy.confinement}
\end{theorem}

\begin{theorem}
A defect paired with its opposite always closes up observably, whatever the block permutation.  No condition on $n$, no fine-tuning.
\begin{align*}
\text{Observable}(M \circ M^{-1})
\end{align*}
\begin{proof}
Unfolding $\text{Observable}$ and composition, the goal is $(M\circ M^{-1}).\text{perm}=1$, i.e.\ $M.\text{perm}\cdot M.\text{perm}^{-1}=1$, which holds for any group element by the cancellation law $g\cdot g^{-1}=1$.
\end{proof}
\lean{Color.BlockMonodromy.pair\_observable}
\end{theorem}

\begin{theorem}
The pair carries zero charge.
\begin{align*}
\text{charge}(M \circ M^{-1}) = 0
\end{align*}
\begin{proof}
By the charge-of-composition law, $\text{charge}(M\circ M^{-1}) = \text{charge}(M) + \text{charge}(M^{-1})$, and by the charge-of-opposite law $\text{charge}(M^{-1})=-\text{charge}(M)$, so the sum is $\text{charge}(M) + (-\text{charge}(M)) = 0$.
\end{proof}
\lean{Color.BlockMonodromy.pair\_charge}
\end{theorem}

\begin{definition}
Bringing together $k$ copies of the same defect is defined recursively:
\begin{align*}
\text{rep}(M, 0) &= \text{triv}(n) \\
\text{rep}(M, k+1) &= M \circ \text{rep}(M, k)
\end{align*}
\lean{Color.BlockMonodromy.rep}
\end{definition}

\begin{proposition}
The permutation of $k$ copies is the $k$-fold power of the single permutation.
\begin{align*}
\text{rep}(M, k).\text{perm} = M.\text{perm}^k
\end{align*}
\begin{proof}
By induction on $k$. At $k=0$, $\text{rep}(M,0)=\text{triv}(n)$ has trivial permutation, and $M.\text{perm}^0=1$, so both sides agree. For the step, assume $\text{rep}(M,m).\text{perm}=M.\text{perm}^m$; then $\text{rep}(M,m+1).\text{perm} = (M\circ\text{rep}(M,m)).\text{perm} = M.\text{perm}\cdot\text{rep}(M,m).\text{perm} = M.\text{perm}\cdot M.\text{perm}^m = M.\text{perm}^{m+1}$, using the composition-permutation law and the inductive hypothesis.
\end{proof}
\lean{Color.BlockMonodromy.rep\_perm}
\end{proposition}

\begin{proposition}
The charge of $k$ copies is $k$ times the single charge.
\begin{align*}
\text{charge}(\text{rep}(M, k)) = (k : \mathbb{Z}) \cdot \text{charge}(M)
\end{align*}
\begin{proof}
By induction on $k$. At $k=0$, $\text{rep}(M,0)=\text{triv}(n)$ has charge $0$ (a sum of zero windings), matching $(0:\mathbb{Z})\cdot\text{charge}(M)=0$. For the step, assume $\text{charge}(\text{rep}(M,m)) = m\cdot\text{charge}(M)$; then by the charge-of-composition law, $\text{charge}(\text{rep}(M,m+1)) = \text{charge}(M) + \text{charge}(\text{rep}(M,m)) = \text{charge}(M) + m\cdot\text{charge}(M) = (m+1)\cdot\text{charge}(M)$, using elementary ring arithmetic on the cast integers.
\end{proof}
\lean{Color.BlockMonodromy.rep\_charge}
\end{proposition}

\begin{theorem}
$k$ copies close up observably exactly when the permutation raised to the $k$-th power equals the identity.
\begin{align*}
\text{If } M.\text{perm}^k = 1, \text{ then } \text{Observable}(\text{rep}(M, k))
\end{align*}
\begin{proof}
By definition $\text{Observable}(\text{rep}(M,k))$ unfolds to
$\text{rep}(M,k).\text{perm}=1$, and $\text{rep}(M,k).\text{perm}=M.\text{perm}^k$
by the representation-power law above, so the goal reduces exactly to the
hypothesis $M.\text{perm}^k=1$.
\end{proof}
\lean{Color.BlockMonodromy.bound\_state\_observable}
\end{theorem}

\begin{theorem}[Bound state size]
The size of a bound state is the multiplicative order of its monodromy permutation.  $k$ copies close up exactly when $\pi^k = 1$.  For a defect that cycles a block of size $n$ once, the smallest consistent number of constituents is $n$: \textbf{baryon number is the block size}, not an extra quantum number.
\begin{align*}
\text{If } \text{ord}(M.\text{perm}) = n, \text{ then } \text{Observable}(\text{rep}(M, n))
\end{align*}
\begin{proof}
Any group element raised to its own order is the identity:
$M.\text{perm}^{\text{ord}(M.\text{perm})}=1$. Substituting $\text{ord}(M.\text{perm})=n$
gives $M.\text{perm}^n=1$, and the previous theorem (bound state observable)
applied with $k=n$ gives $\text{Observable}(\text{rep}(M,n))$.
\end{proof}
\lean{Color.BlockMonodromy.bound\_state\_size}
\end{theorem}

\begin{theorem}[No partial hadrons]
Fewer constituents than the block size never closes up observably.  With $n = 3$: a single constituent is unobservable, two are unobservable, three are observable.  The absence of two-quark hadrons is not an extra rule but a mathematical consequence.
\begin{align*}
\text{If } \text{ord}(M.\text{perm}) = n, \, 0 < k < n, \text{ then } \neg \text{Observable}(\text{rep}(M, k))
\end{align*}
\begin{proof}
Suppose for contradiction that $\text{Observable}(\text{rep}(M,k))$ holds, i.e.
$M.\text{perm}^k=1$ (using $\text{rep}(M,k).\text{perm}=M.\text{perm}^k$ and
$0<k$). Then the order of $M.\text{perm}$ must divide, and in particular be
$\leq k$: $\text{ord}(M.\text{perm}) \leq k$. Substituting the hypothesis
$\text{ord}(M.\text{perm})=n$ gives $n \leq k$, contradicting $k<n$.
\end{proof}
\lean{Color.BlockMonodromy.no\_partial\_state}
\end{theorem}

\begin{theorem}[Charge in $n$ths]
The bound state's charge is $n$ times a constituent's.  Therefore if observable charge is normalised to integers, constituent charges lie in $(1/n)\mathbb{Z}$.
\begin{align*}
\text{charge}(\text{rep}(M, n)) = (n : \mathbb{Z}) \cdot \text{charge}(M)
\end{align*}
\begin{proof}
This is exactly the representation-charge law ($\text{rep}(M,k).\text{charge} = k\cdot M.\text{charge}$) instantiated at $k=n$.
\end{proof}
\lean{Color.BlockMonodromy.constituent\_charge\_fraction}
\end{theorem}

\begin{theorem}
Stated as a fraction: the constituent carries $1/n$ of the observable charge.
\begin{align*}
(M.\text{charge} : \mathbb{Q}) = \frac{1}{n} \cdot (\text{rep}(M, n).\text{charge} : \mathbb{Q})
\end{align*}
when $0 < n$.
\begin{proof}
By the previous theorem, $\text{rep}(M,n).\text{charge} = n\cdot M.\text{charge}$
(cast into $\mathbb{Q}$). Since $n>0$, $(n:\mathbb{Q})\neq0$, so dividing both
sides by $n$ (equivalently, cancelling $n^{-1}\cdot n=1$) gives
$M.\text{charge} = \tfrac1n\cdot\text{rep}(M,n).\text{charge}$.
\end{proof}
\lean{Color.BlockMonodromy.constituent\_charge\_rat}
\end{theorem}

\begin{theorem}
An observable bound state of $n$ identical constituents has charge divisible by $n$.
\begin{align*}
n \mid \text{charge}(\text{rep}(M, n))
\end{align*}
\begin{proof}
Immediate from the charge-in-$n$ths theorem: $\text{rep}(M,n).\text{charge} = n\cdot M.\text{charge}$ exhibits $M.\text{charge}$ itself as the witnessing quotient.
\end{proof}
\lean{Color.BlockMonodromy.bound\_charge\_divisible}
\end{theorem}

\subsection{ColorAudit}

The identification of the degenerate block from the scale pattern with the colour block fails by a factor of $3/2$ in baryon number and is excluded by the charge quantum — halves are not thirds.  This file records the refutation.  The mathematics of \texttt{Color.lean} is untouched; it always said $n$ was an input.  What is excluded is the route that would derive $n$ from the spatial scale pattern.  That route gives $n = 2$ and is refuted.

\begin{definition}
The size of the degenerate block in a uniaxial pattern in three spatial dimensions is two (one distinguished direction, the other two are not).
\begin{align*}
\text{uniaxialBlockSize} = 2
\end{align*}
\lean{ColorAudit.uniaxialBlockSize}
\end{definition}

\begin{definition}
The number of constituents \texttt{Color.lean} predicts for a degenerate block of size $n$.
\begin{align*}
\text{boundStateSize}(n) = n
\end{align*}
\lean{ColorAudit.boundStateSize}
\end{definition}

\begin{definition}
The charge quantum \texttt{Color.lean} predicts for a block of size $n$.
\begin{align*}
\text{chargeQuantum}(n) = \frac{1}{n}
\end{align*}
\lean{ColorAudit.chargeQuantum}
\end{definition}

\begin{definition}
The observed baryon number in nature.
\begin{align*}
\text{observedBaryonNumber} = 3
\end{align*}
\lean{ColorAudit.observedBaryonNumber}
\end{definition}

\begin{definition}
The observed quark charge quantum.
\begin{align*}
\text{observedChargeQuantum} = \frac{1}{3}
\end{align*}
\lean{ColorAudit.observedChargeQuantum}
\end{definition}

\begin{theorem}
Under the identification, bound states would have two constituents.
\begin{align*}
\text{boundStateSize}(\text{uniaxialBlockSize}) = 2
\end{align*}
\begin{proof}
Immediate from the definitions: $\text{boundStateSize}(n)=n$ and $\text{uniaxialBlockSize}=2$, so $\text{boundStateSize}(\text{uniaxialBlockSize})=2$ by unfolding both.
\end{proof}
\lean{ColorAudit.identification\_predicts\_two}
\end{theorem}

\begin{theorem}
And charges would come in halves.
\begin{align*}
\text{chargeQuantum}(\text{uniaxialBlockSize}) = \frac{1}{2}
\end{align*}
\begin{proof}
Unfold $\text{chargeQuantum}(n)=1/n$ and $\text{uniaxialBlockSize}=2$ to get $1/2=1/2$, a numerical identity.
\end{proof}
\lean{ColorAudit.identification\_predicts\_halves}
\end{theorem}

\begin{theorem}
The identification is excluded by baryon number.
\begin{align*}
\text{boundStateSize}(\text{uniaxialBlockSize}) \neq \text{observedBaryonNumber}
\end{align*}
\begin{proof}
Unfolding the definitions reduces the claim to $2 \neq 3$, which holds by arithmetic on natural numbers.
\end{proof}
\lean{ColorAudit.identification\_fails\_baryon}
\end{theorem}

\begin{theorem}
And by the charge quantum.  Halves are not thirds, and not nearly.
\begin{align*}
\text{chargeQuantum}(\text{uniaxialBlockSize}) \neq \text{observedChargeQuantum}
\end{align*}
\begin{proof}
Unfolding the definitions reduces the claim to $1/2 \neq 1/3$, a numerical fact about rationals.
\end{proof}
\lean{ColorAudit.identification\_fails\_charge}
\end{theorem}

\begin{theorem}
The spatial degeneracy block cannot be the colour block.
\begin{align*}
\text{boundStateSize}(\text{uniaxialBlockSize}) \neq \text{observedBaryonNumber} \,\wedge\, \text{chargeQuantum}(\text{uniaxialBlockSize}) \neq \text{observedChargeQuantum}
\end{align*}
\begin{proof}
The conjunction of the two preceding refutations: pair the baryon-number failure with the charge-quantum failure.
\end{proof}
\lean{ColorAudit.spatial\_block\_is\_not\_colour\_block}
\end{theorem}

\begin{theorem}
Only $n = 3$ matches the observed baryon number.
\begin{align*}
\text{If } \text{boundStateSize}(n) = \text{observedBaryonNumber}, \text{ then } n = 3
\end{align*}
\begin{proof}
Since $\text{boundStateSize}(n)=n$ and $\text{observedBaryonNumber}=3$ by definition, the hypothesis $\text{boundStateSize}(n)=\text{observedBaryonNumber}$ literally is the equation $n=3$.
\end{proof}
\lean{ColorAudit.three\_is\_forced\_by\_data}
\end{theorem}

\begin{theorem}
And $n = 3$ does reproduce the observed charge quantum, so the two observations are consistent with a single block size — just not the spatial one.
\begin{align*}
\text{chargeQuantum}(3) = \text{observedChargeQuantum}
\end{align*}
\begin{proof}
Unfold both sides: $\text{chargeQuantum}(3)=1/3$ and $\text{observedChargeQuantum}=1/3$ by definition, so they are equal.
\end{proof}
\lean{ColorAudit.three\_matches\_charge}
\end{theorem}

\begin{theorem}
The gap, stated as a number: the geometry offers two where the data requires three.
\begin{align*}
\text{uniaxialBlockSize} = 2 \,\wedge\, \text{observedBaryonNumber} = 3 \,\wedge\, \text{uniaxialBlockSize} \neq \text{observedBaryonNumber}
\end{align*}
\begin{proof}
The first two conjuncts unfold directly from the definitions $\text{uniaxialBlockSize}=2$ and $\text{observedBaryonNumber}=3$; the third is the numerical fact $2\neq3$.
\end{proof}
\lean{ColorAudit.geometry\_offers\_two\_data\_needs\_three}
\end{theorem}

\subsection{Codimension}

A loop in three-dimensional space does not enclose a point; it encloses a line.  What a loop detects through the scale (which is circle-valued in this file) is a locus of codimension two.  So a circle-valued scale in three dimensions produces strings, not point particles.  The scale is one number so no second scale offers a way to detect anything else.  Every theorem below is true of the circle-valued case and stands as a record of the error corrected in \texttt{Axis.lean}, where the directional scale of axiom A4 is shown to restore point defects.

\begin{definition}
The dimension of the locus a loop can wind around in $d$ spatial dimensions.  A loop detects codimension two; nothing about a circle-valued scale offers a way to detect anything else.
\begin{align*}
\text{defectDim}(d) = d - 2
\end{align*}
\lean{Codimension.defectDim}
\end{definition}

\begin{proposition}
In two spatial dimensions, winding defects are zero-dimensional (pointlike).
\begin{align*}
\text{defectDim}(2) = 0
\end{align*}
\begin{proof}
Immediate from the definition $\text{defectDim}(d)=d-2$: substituting $d=2$ gives $2-2=0$.
\end{proof}
\lean{Codimension.defectDim\_two}
\end{proposition}

\begin{proposition}
In three spatial dimensions, winding defects are one-dimensional (strings).
\begin{align*}
\text{defectDim}(3) = 1
\end{align*}
\begin{proof}
Immediate from the definition: $\text{defectDim}(3) = 3-2 = 1$.
\end{proof}
\lean{Codimension.defectDim\_three}
\end{proposition}

\begin{proposition}
In four spatial dimensions, defects are two-dimensional.
\begin{align*}
\text{defectDim}(4) = 2
\end{align*}
\begin{proof}
Immediate from the definition: $\text{defectDim}(4) = 4-2 = 2$.
\end{proof}
\lean{Codimension.defectDim\_four}
\end{proposition}

\begin{theorem}
Point-like winding defects require exactly two spatial dimensions.  The picture of a point with a loop around it is available only in $d = 2$.
\begin{align*}
\text{If } 2 \le d, \text{ then } (\text{defectDim}(d) = 0 \Leftrightarrow d = 2)
\end{align*}
\begin{proof}
Both directions use $d\geq2$. ($\Rightarrow$) Unfolding $\text{defectDim}(d)=d-2$ (truncated natural subtraction), $d-2=0$ together with $2\leq d$ forces $d=2$ by elementary arithmetic on $\mathbb{N}$. ($\Leftarrow$) Substituting $d=2$ gives $\text{defectDim}(2)=2-2=0$ by the base computation above.
\end{proof}
\lean{Codimension.pointlike\_iff\_two\_dim}
\end{theorem}

\begin{theorem}
In three spatial dimensions a winding defect is a string — a genuinely one-dimensional object, not a point with internal structure.
\begin{align*}
\text{defectDim}(3) = 1
\end{align*}
\begin{proof}
Immediate from the definition, as above: $\text{defectDim}(3)=3-2=1$.
\end{proof}
\lean{Codimension.three\_dim\_gives\_strings}
\end{theorem}

\begin{definition}
Spacetime dimension: spatial dimension plus the single drift direction that the scale supplies.
\begin{align*}
\text{spacetimeDim}(d) = d + 1
\end{align*}
\lean{Codimension.spacetimeDim}
\end{definition}

\begin{theorem}
If the fundamental excitations are point-like, the world is $2+1$.
\begin{align*}
\text{If } 2 \le d \text{ and } \text{defectDim}(d) = 0, \text{ then } \text{spacetimeDim}(d) = 3
\end{align*}
\begin{proof}
By the point-like-iff-two-dimensional theorem, $\text{defectDim}(d)=0$ together with $2\leq d$ gives $d=2$. Substituting into the definition $\text{spacetimeDim}(d)=d+1$ gives $\text{spacetimeDim}(2)=3$.
\end{proof}
\lean{Codimension.pointlike\_forces\_three}
\end{theorem}

\begin{theorem}
In a $3+1$ world the excitations are one-dimensional.  The framework does not get to have both point-like defects and our spacetime.
\begin{align*}
\text{spacetimeDim}(3) = 4 \,\wedge\, \text{defectDim}(3) = 1
\end{align*}
\begin{proof}
Both conjuncts are direct computations from the definitions: $\text{spacetimeDim}(3)=3+1=4$ and $\text{defectDim}(3)=3-2=1$.
\end{proof}
\lean{Codimension.our\_world\_gives\_strings}
\end{theorem}

\begin{theorem}
The two possibilities are mutually exclusive: no spatial dimension gives both point-like defects and a $3+1$ spacetime.
\begin{align*}
\text{If } 2 \le d \text{ and } \text{spacetimeDim}(d) = 4, \text{ then } \text{defectDim}(d) \neq 0
\end{align*}
\begin{proof}
Suppose for contradiction that $\text{defectDim}(d)=0$. By the point-like-iff-two-dimensional theorem (using $2\leq d$) this gives $d=2$. Substituting into $\text{spacetimeDim}(d)=d+1$ gives $\text{spacetimeDim}(d)=3$, contradicting the hypothesis $\text{spacetimeDim}(d)=4$ (since $3\neq4$).
\end{proof}
\lean{Codimension.no\_pointlike\_in\_four}
\end{theorem}

\begin{definition}
The mass of a one-dimensional defect is tension times length.
\begin{align*}
\text{stringMass}(\text{tension}, \text{length}) = \text{tension} \times \text{length}
\end{align*}
\lean{Codimension.stringMass}
\end{definition}

\begin{theorem}
Winding does not determine mass.  Two defects of the same winding but different length have different mass.  So no function of the winding alone can be the mass law — independently confirming \texttt{MassAudit.mass\_law\_underdetermined} and explaining its cause: the defect is extended and topology never measured the extension.
\begin{align*}
\text{If } 0 < \text{tension} \text{ and } L_1 \neq L_2, \text{ then } \text{stringMass}(\text{tension}, L_1) \neq \text{stringMass}(\text{tension}, L_2)
\end{align*}
\begin{proof}
Suppose for contradiction $\text{stringMass}(\text{tension},L_1)=\text{stringMass}(\text{tension},L_2)$, i.e.\ $\text{tension}\cdot L_1 = \text{tension}\cdot L_2$ by definition of $\text{stringMass}$. Since $\text{tension}>0$, it is nonzero, so left-cancellation gives $L_1=L_2$, contradicting the hypothesis $L_1\neq L_2$.
\end{proof}
\lean{Codimension.string\_mass\_not\_topological}
\end{theorem}

\begin{theorem}
Mass grows strictly with length at fixed tension, so the spectrum of a string defect is continuous in length — not a tower indexed by an integer.
\begin{align*}
\text{If } 0 < \text{tension} \text{ and } L_1 < L_2, \text{ then } \text{stringMass}(\text{tension}, L_1) < \text{stringMass}(\text{tension}, L_2)
\end{align*}
\begin{proof}
Unfolding $\text{stringMass}(\text{tension},L)=\text{tension}\cdot L$, the claim is $\text{tension}\cdot L_1 < \text{tension}\cdot L_2$, which follows from $L_1<L_2$ by multiplying a strict inequality on the left by the positive constant $\text{tension}$.
\end{proof}
\lean{Codimension.stringMass\_strictMono}
\end{theorem}

\begin{proposition}
A closed string of zero length is massless.
\begin{align*}
\text{stringMass}(\text{tension}, 0) = 0
\end{align*}
\begin{proof}
Immediate from the definition: $\text{stringMass}(\text{tension},0) = \text{tension}\times0 = 0$.
\end{proof}
\lean{Codimension.stringMass\_zero}
\end{proposition}

\begin{theorem}
For any spatial dimension at least two, either the defects are extended or the world is $2+1$.
\begin{align*}
\text{If } 2 \le d, \text{ then } (\text{defectDim}(d) \neq 0 \,\vee\, \text{spacetimeDim}(d) = 3)
\end{align*}
\begin{proof}
Case split on whether $d=2$ or $2<d$ (the trichotomy available from $2\leq d$). If $d=2$, take the right disjunct: $\text{spacetimeDim}(2)=2+1=3$. If $2<d$, take the left disjunct: unfolding $\text{defectDim}(d)=d-2$ (truncated subtraction), $2<d$ gives $d-2\neq0$ by elementary arithmetic on $\mathbb{N}$.
\end{proof}
\lean{Codimension.extended\_or\_low\_dimensional}
\end{theorem}

\subsection{Expressive}

The framework's vocabulary is three operations: comparing scales, failing to commute, and counting.  This file proves these really are the whole vocabulary and that they are mutually independent.  The scale channel expresses differences and nothing else, rotation says nothing about scale and scale says nothing about rotation, so the two channels are disjoint.  A cross-sector test would need them to constrain each other, and the axioms guarantee they do not.  These theorems hold of the scalar scale sector; the directional scale of A4 couples the channels, as shown in \texttt{Coupling.lean}.

\begin{definition}
A quantity is \textbf{observable} when a global change of the fiducial (an additive shift of the scale field everywhere) does not move it.  This is axiom A5 stated for an arbitrary function of the scale field.
\begin{align*}
\text{FiducialInvariant}(f) :\Leftrightarrow \forall (\sigma : X \to \mathbb{R})(c : \mathbb{R}), \, f(\lambda x. \sigma x + c) = f(\sigma)
\end{align*}
\lean{Expressive.FiducialInvariant}
\end{definition}

\begin{definition}
The difference pattern of a scale field relative to a base point $x_0$ is the residual information after accounting for the global shift.
\begin{align*}
\text{diffPattern}(x_0, \sigma) = \lambda x. \sigma x - \sigma x_0
\end{align*}
\lean{Expressive.diffPattern}
\end{definition}

\begin{proposition}
Adding a constant to the scale field does not change the difference pattern.
\begin{align*}
\text{diffPattern}(x_0, \lambda x. \sigma x + c) = \text{diffPattern}(x_0, \sigma)
\end{align*}
\begin{proof}
Evaluate both sides at an arbitrary $x$: the left side is $(\sigma x + c) - (\sigma x_0 + c)$, which simplifies by cancelling the additive constant $c$ to $\sigma x - \sigma x_0$, the right side. Since the two functions agree pointwise, they are equal.
\end{proof}
\lean{Expressive.diffPattern\_shift}
\end{proposition}

\begin{theorem}
Every observable is a function of differences.  Given invariance, evaluating on the difference pattern gives the same answer as evaluating on the field itself.  So no observable can depend on anything beyond differences.
\begin{align*}
\text{If } \text{FiducialInvariant}(f), \text{ then } f(\sigma) = f(\text{diffPattern}(x_0, \sigma))
\end{align*}
\begin{proof}
Apply the invariance hypothesis to the field $\text{diffPattern}(x_0,\sigma)$ shifted by the constant $c=\sigma(x_0)$: $f(\text{diffPattern}(x_0,\sigma)) = f(\lambda x.\,\text{diffPattern}(x_0,\sigma)(x) + \sigma(x_0))$. Since $\text{diffPattern}(x_0,\sigma)(x)+\sigma(x_0) = (\sigma x - \sigma x_0) + \sigma x_0 = \sigma x$ for every $x$, the shifted field is $\sigma$ itself, so this equals $f(\sigma)$. Reading the chain backwards gives $f(\sigma)=f(\text{diffPattern}(x_0,\sigma))$.
\end{proof}
\lean{Expressive.invariant\_factors}
\end{theorem}

\begin{theorem}
Every function of differences is an observable.
\begin{align*}
\text{FiducialInvariant}(\lambda \sigma. g(\text{diffPattern}(x_0, \sigma)))
\end{align*}
\begin{proof}
Fix $\sigma$ and $c$; the goal is $g(\text{diffPattern}(x_0,\lambda x.\sigma x+c)) = g(\text{diffPattern}(x_0,\sigma))$. By the shift-invariance of the difference pattern proved above, $\text{diffPattern}(x_0,\lambda x.\sigma x+c)=\text{diffPattern}(x_0,\sigma)$, so applying $g$ to both sides gives the goal.
\end{proof}
\lean{Expressive.diffPattern\_invariant}
\end{theorem}

\begin{theorem}
The scale channel expresses differences and nothing else.  Observables are exactly the functions of the difference pattern — the inclusion holds in both directions.  This is the precise expressive limit that A5 imposes.
\begin{align*}
\text{FiducialInvariant}(f) \Leftrightarrow \exists g, \forall \sigma, \, f(\sigma) = g(\text{diffPattern}(x_0, \sigma))
\end{align*}
\begin{proof}
($\Rightarrow$) Given $\text{FiducialInvariant}(f)$, take $g=f$ itself; the required identity $f(\sigma)=f(\text{diffPattern}(x_0,\sigma))$ is exactly the invariant-factors theorem above. ($\Leftarrow$) Given $g$ with $f(\sigma)=g(\text{diffPattern}(x_0,\sigma))$ for all $\sigma$, fix $\sigma,c$; then $f(\lambda x.\sigma x+c) = g(\text{diffPattern}(x_0,\lambda x.\sigma x+c)) = g(\text{diffPattern}(x_0,\sigma)) = f(\sigma)$, using the hypothesis twice and the shift-invariance of the difference pattern in between.
\end{proof}
\lean{Expressive.invariant\_iff\_diffPattern}
\end{theorem}

\begin{theorem}
Absolute scale is not expressible.  Reading the scale at a point is not invariant, so it is not an observable of this framework at all — not merely unmeasured, but outside the vocabulary.
\begin{align*}
\neg \text{FiducialInvariant}(\lambda \sigma. \sigma(x))
\end{align*}
\begin{proof}
Suppose for contradiction that evaluation at $x$ were invariant. Apply the invariance to the zero field $\sigma\equiv0$ and shift $c=1$: invariance demands $(\lambda x'. 0+1)(x) = 0(x)$, i.e.\ $1=0$, which is false. Hence evaluation is not invariant.
\end{proof}
\lean{Expressive.eval\_not\_invariant}
\end{theorem}

\begin{theorem}
The scale of a comparison says nothing about its rotation.  Two comparisons with the same scale differ by an arbitrary element of the commutator subgroup, so fixing the scale leaves rotation entirely free.
\begin{align*}
\text{If } k \in \text{commutator}(\Gamma), \text{ then } \text{scale}(k \cdot g) = \text{scale}(g)
\end{align*}
\begin{proof}
Since $\text{scale}$ is a homomorphism, $\text{scale}(k\cdot g) = \text{scale}(k)\cdot\text{scale}(g)$. Because $k$ lies in the commutator subgroup, it lies in the kernel of the abelianization map, so $\text{scale}(k)=1$. Hence $\text{scale}(k\cdot g) = 1\cdot\text{scale}(g) = \text{scale}(g)$.
\end{proof}
\lean{Expressive.scale\_says\_nothing\_about\_rotation}
\end{theorem}

\begin{theorem}
The rotation says nothing about the scale.  Every commutator has trivial scale, so knowing the rotation part determines nothing about the scale part.
\begin{align*}
\text{scale}(g \cdot h \cdot g^{-1} \cdot h^{-1}) = 1
\end{align*}
\begin{proof}
The element $g\cdot h\cdot g^{-1}\cdot h^{-1}$ is a commutator, hence lies in the kernel of the abelianization homomorphism $\text{scale}$; every commutator maps to the identity, giving $\text{scale}(g\cdot h\cdot g^{-1}\cdot h^{-1})=1$.
\end{proof}
\lean{Expressive.rotation\_says\_nothing\_about\_scale}
\end{theorem}

\begin{theorem}
The channels are independent.  Fixing either one constrains the other not at all.  This is by construction — the scale channel is a quotient and the rotation channel is its kernel — so no amount of further work inside the present axioms can make one test the other.
\begin{align*}
(\forall k \in \text{commutator}(\Gamma), \, \text{scale}(k \cdot g) = \text{scale}(g)) \,\wedge\, (\forall a, b : \Gamma, \, \text{scale}(a \cdot b \cdot a^{-1} \cdot b^{-1}) = 1)
\end{align*}
\begin{proof}
The two conjuncts are exactly the two preceding theorems: the first is scale-says-nothing-about-rotation applied for each $k$, and the second is rotation-says-nothing-about-scale applied for each pair $a,b$.
\end{proof}
\lean{Expressive.channels\_independent}
\end{theorem}

\begin{theorem}
A comparison is recovered from its scale only up to a rotation, and from its rotation only up to a scale: neither channel alone is complete, and together they are just the original comparison again.
\begin{align*}
\text{If } \text{scale}(g) = \text{scale}(h), \text{ then } g \cdot h^{-1} \in \text{commutator}(\Gamma)
\end{align*}
\begin{proof}
This is the standard characterization of the abelianization map: two elements have equal scale exactly when they differ by an element of the commutator subgroup (the kernel of the abelianization homomorphism). Applying this characterization to the hypothesis $\text{scale}(g)=\text{scale}(h)$ gives $g\cdot h^{-1}\in\text{commutator}(\Gamma)$.
\end{proof}
\lean{Expressive.neither\_channel\_complete}
\end{theorem}

\begin{theorem}
Within the present axioms, no function of the scale channel can determine the rotation channel, because the scale is blind to exactly the subgroup the rotation lives in.
\begin{align*}
\forall (F : \text{Abelianization}(\Gamma) \to \Gamma)(g : \Gamma)(k : \Gamma), \, k \in \text{commutator}(\Gamma) \Rightarrow \text{scale}(k \cdot g) = \text{scale}(g)
\end{align*}
\begin{proof}
The statement does not depend on $F$ at all — it is scale-says-nothing-about-rotation, restated with an unused universally quantified function $F$ to emphasise that no such $F$ could recover the rotation channel from the scale channel; the proof is that theorem applied to $g$ and $k$.
\end{proof}
\lean{Expressive.no\_channel\_coupling}
\end{theorem}

\subsection{CrossCheck}

A threshold density $\rho$ appears in two places: in the mass spectrum as the average rate of thresholds per $e$-fold, and in the running of couplings as the rate-controlling density.  Since the framework names them the same, the natural move is to measure one from the spectrum and one from running and check they agree — a genuine cross-sector test.  That test was run.  It is not a test, because the two quantities are not the same: one is a raw count of thresholds, the other is sector-restricted and signed.  The proposal was an over-reach and this file records why.

\begin{definition}
Threshold density from all twelve Standard Model mass scales, per $e$-fold.
\begin{align*}
\rho_{\text{all}} = \frac{11}{12.731}
\end{align*}
\lean{CrossCheck.rhoAll}
\end{definition}

\begin{definition}
Threshold density from the six quarks alone — the structures a strong coupling responds to.
\begin{align*}
\rho_{\text{quark}} = \frac{5}{11.271}
\end{align*}
\lean{CrossCheck.rhoQuark}
\end{definition}

\begin{proposition}
$\rho_{\text{all}}$ lies between $0.86$ and $0.87$.
\begin{align*}
0.86 < \rho_{\text{all}} < 0.87
\end{align*}
\begin{proof}
Substitute $\rho_{\text{all}}=11/12.731$ and evaluate: $11/12.731 \approx 0.8640$, which lies strictly between $0.86$ and $0.87$ by direct numerical computation.
\end{proof}
\lean{CrossCheck.rhoAll\_value}
\end{proposition}

\begin{proposition}
$\rho_{\text{quark}}$ lies between $0.44$ and $0.45$.
\begin{align*}
0.44 < \rho_{\text{quark}} < 0.45
\end{align*}
\begin{proof}
Substitute $\rho_{\text{quark}}=5/11.271$ and evaluate: $5/11.271 \approx 0.4436$, which lies strictly between $0.44$ and $0.45$ by direct numerical computation.
\end{proof}
\lean{CrossCheck.rhoQuark\_value}
\end{proposition}

\begin{theorem}
The two densities disagree.
\begin{align*}
\rho_{\text{all}} \neq \rho_{\text{quark}}
\end{align*}
\begin{proof}
Suppose for contradiction $\rho_{\text{all}}=\rho_{\text{quark}}$. By the two bounding propositions above, $0.86 < \rho_{\text{all}}$ and $\rho_{\text{quark}} < 0.45$; substituting the assumed equality gives $0.86 < \rho_{\text{quark}} < 0.45$, an immediate numerical contradiction since $0.86 > 0.45$.
\end{proof}
\lean{CrossCheck.densities\_differ}
\end{theorem}

\begin{theorem}
By roughly a factor of two — which is a fact about how much of the Standard Model happens to be coloured, not a prediction of anything.
\begin{align*}
1.9 < \frac{\rho_{\text{all}}}{\rho_{\text{quark}}} < 2.0
\end{align*}
\begin{proof}
By the two bounding propositions, $\rho_{\text{quark}}>0.44>0$, so multiplying through by $\rho_{\text{quark}}$ preserves the direction of the inequalities. The lower bound $1.9<\rho_{\text{all}}/\rho_{\text{quark}}$ is equivalent, after clearing the denominator, to $1.9\cdot\rho_{\text{quark}}<\rho_{\text{all}}$, which follows from $\rho_{\text{quark}}<0.45$ and $\rho_{\text{all}}>0.86$ by a numerical bound ($1.9\times0.45=0.855<0.86$). The upper bound $\rho_{\text{all}}/\rho_{\text{quark}}<2.0$ is equivalent to $\rho_{\text{all}}<2.0\cdot\rho_{\text{quark}}$, which follows from $\rho_{\text{all}}<0.87$ and $\rho_{\text{quark}}>0.44$ ($2.0\times0.44=0.88>0.87$).
\end{proof}
\lean{CrossCheck.density\_ratio\_near\_two}
\end{theorem}

\begin{theorem}
Within a sector, the relation has no content.  One equation and one unknown: given measurements of the beta drift and the threshold density this defines the weight $\kappa$, and substituting back returns the input.  Nothing has been tested.
\begin{align*}
\text{If } dN/dt \neq 0, \text{ then } (db/dt / dN/dt) \cdot dN/dt = db/dt
\end{align*}
\begin{proof}
Since $dN/dt\neq0$, division by it and multiplication by it are mutually cancelling field operations: $(db/dt \,/\, dN/dt)\cdot dN/dt = db/dt \cdot (dN/dt)^{-1}\cdot dN/dt = db/dt\cdot 1 = db/dt$.
\end{proof}
\lean{CrossCheck.sector\_relation\_is\_identity}
\end{theorem}

\begin{theorem}
Carried out for QCD between the charm and top thresholds: two flavours appear, the beta coefficient drifts by $2/3$ per flavour, and the extracted weight is $2/3$ — the input written backwards.
\begin{align*}
\text{If } dN/dt \neq 0, \text{ then } \left(\frac{2}{3} \cdot dN/dt\right) / dN/dt = \frac{2}{3}
\end{align*}
\begin{proof}
Since $dN/dt\neq0$, it is invertible in the field of reals, so $(\tfrac23\cdot dN/dt)/dN/dt = \tfrac23\cdot dN/dt\cdot(dN/dt)^{-1} = \tfrac23\cdot1 = \tfrac23$.
\end{proof}
\lean{CrossCheck.qcd\_extraction\_returns\_input}
\end{theorem}

\begin{theorem}
What a genuine test would require: two independent determinations of the same weight, so that the second can check the first.  The framework supplies only one, because it has only one coupling-like quantity per sector and no coupling that responds to every structure.
\begin{align*}
\text{If } dN/dt \neq 0 \text{ and } \kappa = db/dt / dN/dt, \text{ then } db/dt = \kappa \cdot dN/dt
\end{align*}
\begin{proof}
Substitute the definition $\kappa = db/dt / dN/dt$ into the goal, reducing it to $db/dt = (db/dt / dN/dt)\cdot dN/dt$, which is exactly the sector-relation-is-identity theorem above (using $dN/dt\neq0$).
\end{proof}
\lean{CrossCheck.no\_second\_determination}
\end{theorem}

\begin{definition}
Sum of the five quark log-gaps.
\begin{align*}
\text{quarkGapSum} = 11.270
\end{align*}
\lean{CrossCheck.quarkGapSum}
\end{definition}

\begin{definition}
Sum of their squares.
\begin{align*}
\text{quarkGapSqSum} = 31.583624
\end{align*}
\lean{CrossCheck.quarkGapSqSum}
\end{definition}

\begin{definition}
Mean of the five quark log-gaps.
\begin{align*}
\text{quarkMeanGap} = \text{quarkGapSum} / 5
\end{align*}
\lean{CrossCheck.quarkMeanGap}
\end{definition}

\begin{definition}
Variance of the five quark log-gaps.
\begin{align*}
\text{quarkVariance} = \text{quarkGapSqSum} / 5 - \text{quarkMeanGap}^2
\end{align*}
\lean{CrossCheck.quarkVariance}
\end{definition}

\begin{definition}
Squared coefficient of variation for the quark sector.
\begin{align*}
\text{quarkCvSq} = \text{quarkVariance} / \text{quarkMeanGap}^2
\end{align*}
\lean{CrossCheck.quarkCvSq}
\end{definition}

\begin{theorem}
$CV^2 = 0.243$ (so $CV = 0.493$) — the twelfth percentile of the constant-rate prediction for five gaps.  Low, and more regular than the prediction typically gives, but not excluded.  What survives is the per-sector prediction: each sector's thresholds should be a constant-rate process.
\begin{align*}
0.24 < \text{quarkCvSq} < 0.25
\end{align*}
\begin{proof}
Unfold the chain of definitions: $\text{quarkMeanGap}=11.270/5=2.254$, $\text{quarkVariance}=31.583624/5 - 2.254^2 = 6.3167248 - 5.080516 = 1.2362088$, and $\text{quarkCvSq}=1.2362088/2.254^2 = 1.2362088/5.080516\approx0.24333$. This lies strictly between $0.24$ and $0.25$ by direct numerical computation.
\end{proof}
\lean{CrossCheck.quarkCvSq\_value}
\end{theorem}

\begin{theorem}
The quark sector is not a geometric tower: its gap variance is positive.
\begin{align*}
\text{quarkCvSq} \neq 0
\end{align*}
\begin{proof}
By the previous theorem $0.24<\text{quarkCvSq}$, so in particular $\text{quarkCvSq}$ is strictly positive and hence cannot equal $0$.
\end{proof}
\lean{CrossCheck.quark\_not\_geometric}
\end{theorem}

\subsection{MassAudit}

The mass law $m_k = m_0 e^{|k|\Delta}$ was presented as a consequence of axiom A3 ("mass is inverse scale") and the scale stepping by $\Delta$ per unit winding.  Testing it against lattice glueball data falsified the constant mass ratio.  But the law was never derived — it was assumed.  The topology establishes only that winding is an integer, additive, conserved, and invariant under orientation reversal.  What the mass law additionally requires — how strained the scale field is by a defect of winding $k$ — is a dynamical question about energetics, which no theorem addresses.  This file exhibits two different mass laws, exponential and quadratic, both fully compatible with every proved property of winding, so the topology does not fix the mass law.

\begin{definition}
The law asserted in earlier work: a geometric tower.
\begin{align*}
\text{massExp}(m_0, \Delta, k) = m_0 \cdot e^{|k|\Delta}
\end{align*}
\lean{MassAudit.massExp}
\end{definition}

\begin{definition}
An equally available alternative: strain energy quadratic in the winding, which is what a naive energetic estimate for a topological defect gives.
\begin{align*}
\text{massQuad}(m_0, c, k) = m_0 + c \cdot k^2
\end{align*}
\lean{MassAudit.massQuad}
\end{definition}

\begin{proposition}
The exponential law satisfies particle–antiparticle mass degeneracy.
\begin{align*}
\text{massExp}(m_0, \Delta, -k) = \text{massExp}(m_0, \Delta, k)
\end{align*}
\begin{proof}
Unfold $\text{massExp}(m_0,\Delta,k) = m_0\cdot e^{|k|\Delta}$. Since $|-k|=|k|$ for any integer $k$, the exponents agree, so $\text{massExp}(m_0,\Delta,-k) = m_0\cdot e^{|-k|\Delta} = m_0\cdot e^{|k|\Delta} = \text{massExp}(m_0,\Delta,k)$.
\end{proof}
\lean{MassAudit.massExp\_neg}
\end{proposition}

\begin{proposition}
The quadratic law also satisfies particle–antiparticle mass degeneracy.
\begin{align*}
\text{massQuad}(m_0, c, -k) = \text{massQuad}(m_0, c, k)
\end{align*}
\begin{proof}
Unfold $\text{massQuad}(m_0,c,k) = m_0+c\cdot k^2$. Since $(-k)^2=k^2$ for any integer $k$ cast to $\mathbb{R}$, $\text{massQuad}(m_0,c,-k) = m_0 + c\cdot(-k)^2 = m_0+c\cdot k^2 = \text{massQuad}(m_0,c,k)$ by ring arithmetic.
\end{proof}
\lean{MassAudit.massQuad\_neg}
\end{proposition}

\begin{theorem}
The exponential law is positive for positive parameters.
\begin{align*}
\text{If } 0 < m_0, \text{ then } 0 < \text{massExp}(m_0, \Delta, k)
\end{align*}
\begin{proof}
Unfold $\text{massExp}(m_0,\Delta,k)=m_0\cdot e^{|k|\Delta}$: the exponential factor $e^{|k|\Delta}$ is always strictly positive, and $m_0>0$ by hypothesis, so the product of two positive reals is positive.
\end{proof}
\lean{MassAudit.massExp\_pos}
\end{theorem}

\begin{theorem}
The quadratic law is also positive for positive parameters.
\begin{align*}
\text{If } 0 < m_0 \text{ and } 0 \le c, \text{ then } 0 < \text{massQuad}(m_0, c, k)
\end{align*}
\begin{proof}
Unfold $\text{massQuad}(m_0,c,k)=m_0+c\cdot k^2$. Since $c\geq0$ and $k^2\geq0$, the product $c\cdot k^2\geq0$; adding this nonnegative quantity to $m_0>0$ preserves positivity: $0<m_0\leq m_0+c\cdot k^2$.
\end{proof}
\lean{MassAudit.massQuad\_pos}
\end{theorem}

\begin{proposition}
Both laws agree at zero winding.
\begin{align*}
\text{massExp}(m_0, \Delta, 0) = m_0 = \text{massQuad}(m_0, c, 0)
\end{align*}
\begin{proof}
For the first equality, unfold $\text{massExp}(m_0,\Delta,0) = m_0\cdot e^{|0|\Delta} = m_0\cdot e^{0} = m_0\cdot1=m_0$, using $|0|=0$ and $e^0=1$. For the second, unfold $\text{massQuad}(m_0,c,0) = m_0+c\cdot0^2 = m_0+0=m_0$.
\end{proof}
\lean{MassAudit.massExp\_zero}
\end{proposition}

\begin{proposition}
Both laws agree at zero winding (quadratic version).
\begin{align*}
\text{massQuad}(m_0, c, 0) = m_0
\end{align*}
\begin{proof}
Immediate from the definition: $\text{massQuad}(m_0,c,0) = m_0 + c\cdot0^2 = m_0+0 = m_0$ by ring arithmetic.
\end{proof}
\lean{MassAudit.massQuad\_zero}
\end{proposition}

\begin{theorem}
The topology does not determine the mass law.  Two laws that are genuinely different — their ratio at winding $1$ differs — both satisfy every property the topology has fixed: positivity, agreement at $k=0$, and invariance under orientation reversal.  Therefore the geometric tower is an additional hypothesis, and the constant mass ratio it predicts is a consequence of that hypothesis rather than of the framework.
\begin{align*}
\exists f, g : \mathbb{Z} \to \mathbb{R}, \quad &(\forall k, f(-k) = f(k)) \,\wedge\, (\forall k, g(-k) = g(k)) \\
&\,\wedge\, (\forall k, 0 < f(k)) \,\wedge\, (\forall k, 0 < g(k)) \\
&\,\wedge\, (f(0) = g(0)) \,\wedge\, (f(1) \neq g(1))
\end{align*}
\begin{proof}
Take $f=\text{massExp}(1,1,\cdot)$ and $g=\text{massQuad}(1,1,\cdot)$. Orientation-reversal symmetry for both is $\text{massExp\_neg}$ and $\text{massQuad\_neg}$ above; positivity of both for all $k$ is $\text{massExp\_pos}$ and $\text{massQuad\_pos}$ (with $m_0=1>0$, $c=1\geq0$). Agreement at $k=0$ is $\text{massExp\_zero}$ and $\text{massQuad\_zero}$, both giving $1$. For the final inequality, compute $g(1)=\text{massQuad}(1,1,1) = 1+1\cdot1^2=2$, while $f(1)=\text{massExp}(1,1,1)=1\cdot e^{1}=e$. Since $e>2.7$ (a standard numerical bound on Euler's constant) while $g(1)=2<2.7$, $f(1)\neq g(1)$.
\end{proof}
\lean{MassAudit.mass\_law\_underdetermined}
\end{theorem}

\begin{theorem}
Both the exponential and quadratic laws reproduce the one mass statement that is derived — the particle–antiparticle degeneracy — so that prediction survives the retraction while the tower does not.
\begin{align*}
\text{massExp}(m_0, \Delta, -k) = \text{massExp}(m_0, \Delta, k) \,\wedge\, \text{massQuad}(m_0, c, -k) = \text{massQuad}(m_0, c, k)
\end{align*}
\begin{proof}
The two conjuncts are exactly $\text{massExp\_neg}$ and $\text{massQuad\_neg}$ proved above, applied to the given parameters.
\end{proof}
\lean{MassAudit.both\_laws\_antiparticle\_degenerate}
\end{theorem}

\begin{theorem}
The quadratic law has a non-constant consecutive ratio, so a spectrum that fails the geometric-tower test is not thereby in conflict with a defect picture at all.
\begin{align*}
\text{massQuad}(1, 1, 1) \cdot \text{massQuad}(1, 1, 1) \neq \text{massQuad}(1, 1, 2) \cdot \text{massQuad}(1, 1, 0)
\end{align*}
\begin{proof}
Unfold and compute each factor: $\text{massQuad}(1,1,1)=1+1^2=2$, $\text{massQuad}(1,1,2)=1+2^2=5$, $\text{massQuad}(1,1,0)=1+0^2=1$. So the left side is $2\cdot2=4$ and the right side is $5\cdot1=5$, and $4\neq5$ by numerical computation.
\end{proof}
\lean{MassAudit.massQuad\_ratio\_not\_constant}
\end{theorem}

\subsection{Audit}

This file is the register — the single authoritative list of what was retracted, corrected, cited, assumed, or proved with no auxiliary hypotheses. Nothing here is new mathematics; it is the standing answer to ``Does this rest on the axioms alone?'' and it is meant to be read before every claim.  The retractions are real: a geometric mass tower (assumed, not derived), a colour identification (wrong by a factor of $3/2$), a cross-sector test (measuring two different quantities), and a codimension analysis (true of a circle-valued scale, not this framework). The corrections are more numerous than the retractions — most claims survived with their scope narrowed. An earlier version of this register was incomplete and omitted A6$^\prime$ (the scale response law) and the assumption that defects are uniform. The register grew under external audit, which is the point of having one.

\begin{definition}
How many times the scalar-for-directional substitution has been made and caught.  Recorded so the count cannot quietly drift.
\begin{align*}
\text{scalarForDirectionalCount} = 5
\end{align*}
\lean{Audit.scalarForDirectionalCount}
\end{definition}

\begin{theorem}
The scalar-for-directional substitution appears exactly five times in the development.
\begin{align*}
\text{scalarForDirectionalCount} = 5
\end{align*}
\begin{proof}
Immediate from the definition $\text{scalarForDirectionalCount}=5$.
\end{proof}
\lean{Audit.scalar\_for\_directional\_count}
\end{theorem}

\begin{definition}
Claims retracted outright.
\begin{align*}
\text{retractedCount} = 4
\end{align*}
\lean{Audit.retractedCount}
\end{definition}

\begin{definition}
Claims corrected in scope, where the theorem survives and the gloss did not.
\begin{align*}
\text{correctedCount} = 6
\end{align*}
\lean{Audit.correctedCount}
\end{definition}

\begin{definition}
External results cited and not proved.
\begin{align*}
\text{citedCount} = 2
\end{align*}
\lean{Audit.citedCount}
\end{definition}

\begin{definition}
Assumptions registered rather than smuggled: A6$^\prime$ (the scale response law), the assumption that defects are uniform in the symmetric space's volume, and A7 itself.
\begin{align*}
\text{assumedCount} = 3
\end{align*}
\lean{Audit.assumedCount}
\end{definition}

\begin{theorem}
The register is not empty.  A development with nothing retracted has not been audited.  The counts are carried so that a reader can weigh the results against what had to be taken back to get them.
\begin{align*}
\text{retractedCount} \neq 0 \,\wedge\, \text{correctedCount} \neq 0 \,\wedge\, \text{citedCount} \neq 0 \,\wedge\, \text{assumedCount} \neq 0
\end{align*}
\begin{proof}
Each conjunct unfolds to a concrete numeral inequality: $\text{retractedCount}=4\neq0$, $\text{correctedCount}=6\neq0$, $\text{citedCount}=2\neq0$, $\text{assumedCount}=3\neq0$, each true by direct computation on natural numbers.
\end{proof}
\lean{Audit.register\_nonempty}
\end{theorem}

\begin{theorem}
Corrections outnumber retractions: more claims survived with their scope narrowed than had to be withdrawn.
\begin{align*}
\text{retractedCount} < \text{correctedCount}
\end{align*}
\begin{proof}
Unfolding both definitions reduces the claim to $4<6$, true by direct computation on natural numbers.
\end{proof}
\lean{Audit.more\_corrected\_than\_retracted}
\end{theorem}

\begin{definition}
Structures carrying hypotheses the framework does not supply — true theorems whose physical reading is conditional.  These are theorems about anything meeting their hypothesis, not statements of what the axioms deliver.
\begin{align*}
\text{unwitnessedStructures} = 6
\end{align*}
\lean{Audit.unwitnessedStructures}
\end{definition}

\begin{theorem}
The register grew under external audit, which is the point of having one.  An audit that finds nothing has not been run adversarially.
\begin{align*}
0 < \text{unwitnessedStructures}
\end{align*}
\begin{proof}
Unfolding the definition reduces the claim to $0<6$, true by direct computation on natural numbers.
\end{proof}
\lean{Audit.register\_grew\_under\_audit}
\end{theorem}

\begin{definition}
Number of theorems put through axiom verification — every theorem in the development, generated from the sources rather than curated.
\begin{align*}
\text{auditedTheorems} = 650
\end{align*}
\lean{Audit.auditedTheorems}
\end{definition}

\begin{definition}
Occurrences of unresolved axioms in that audit.
\begin{align*}
\text{sorryAxCount} = 0
\end{align*}
\lean{Audit.sorryAxCount}
\end{definition}

\begin{theorem}
Nothing in the development rests on a gap, as a number rather than a mood.  Every audited theorem reduces to `propext`, `Classical.choice` and `Quot.sound` and to nothing else.
\begin{align*}
\text{sorryAxCount} = 0 \,\wedge\, 0 < \text{auditedTheorems}
\end{align*}
\begin{proof}
The first conjunct is the definition $\text{sorryAxCount}=0$ itself. The second unfolds $\text{auditedTheorems}=650$ to the numerical fact $0<650$, true by direct computation.
\end{proof}
\lean{Audit.clean\_count}
\end{theorem}
